% GENERATED FILE. Edit canonical lessons, then run npm run generate:books.
% canonical-source-hash: fnv1a64:fccc089d0692bf6b
% canonical-lessons: TA-C346-uccakattam, TA-C346-pirkalattil, TA-C346-enrenrum, TA-C346-ivvaru, TA-C346-viti-dawn, TA-R346-review-340-346, TA-R346-review-final-346

\chapter{How It Ended}
\label{ch:ta-how-it-ended}
\hlchaptermodality{346}

\begin{chapteropening}
\textbf{By the end of this chapter:} \emph{I can bring a story to its end.}

Complete the last lesson of chapter 346: I can bring a story to its end.
\end{chapteropening}

\section[\texorpdfstring{\ta{உச்சகட்டம்} (uccakaṭṭam)}{உச்சகட்டம் (uccakaṭṭam)}]{\ta{உச்சகட்டம்} (uccakaṭṭam) — a climax; a peak}

\label{lesson:TA-C346-uccakattam}



\begin{quote}
Before the new one: say the Tamil for a setting, then the Tamil for to be staged.
\end{quote}

\subsection*{You'll want to know: \ta{உச்சகட்டம்}}
\textbf{\ta{உச்சகட்டம்}} — \emph{uccakaṭṭam} — \textquotedblleft{}a climax; a peak\textquotedblright{}.

\begin{quote}
\ta{கதையின்} \ta{உச்சகட்டத்தில்} \ta{எல்லோரும்} \ta{அமைதியாக} \ta{இருந்தார்கள்}.
\end{quote}

\begin{quote}
\emph{kataiyiṉ uccakaṭṭattil ellōrum amaitiyāka iruntārkaḷ.}
\end{quote}

\begin{quote}
At the climax of the story everyone was silent.
\end{quote}

\ta{உச்சம்}, peak + \ta{கட்டம்}, stage.

\begin{cousinweb}
\textbf{\ta{உச்சகட்டம்}} (\emph{uccakaṭṭam}) is built from \ta{உச்சம்}, \textquotedblleft{}the highest point\textquotedblright{}, borrowed from Sanskrit \emph{ucca}, \textquotedblleft{}tall, lofty\textquotedblright{} (CDIAL 1634), with \ta{கட்டம்}, \textquotedblleft{}a stage\textquotedblright{}.
\end{cousinweb}

\begin{grammarlens}[title={What you've built}]
One more word for telling what happened.
\end{grammarlens}

\subsection*{Your turn}
\begin{itemize}
\raggedright
  \item \emph{Say it:} \emph{uccakaṭṭam}
  \item \emph{Say it:} \emph{uccakaṭṭam}, once more
  \item \emph{Say it:} \emph{kataiyiṉ uccakaṭṭattil ellōrum amaitiyāka iruntārkaḷ.} — At the climax of the story everyone was silent.
  \item \emph{Recall:} say the Tamil for a setting, then the Tamil for to be staged, then say \emph{uccakaṭṭam} again
\end{itemize}

\begin{tcolorbox}[breakable,colback=teal!4,colframe=teal!35!black,title={Before you move on}]
What does \ta{உச்சகட்டம்} mean? (\textquotedblleft{}A climax; a peak\textquotedblright{}.) Say it once more.
\end{tcolorbox}
\section[\texorpdfstring{\ta{பிற்காலத்தில்} (piṟkālattil)}{பிற்காலத்தில் (piṟkālattil)}]{\ta{பிற்காலத்தில்} (piṟkālattil) — in later times}

\label{lesson:TA-C346-pirkalattil}



\begin{quote}
Before the new one: say the Tamil for to be staged, then the Tamil for a climax.
\end{quote}

\subsection*{You'll want to know: \ta{பிற்காலத்தில்}}
\textbf{\ta{பிற்காலத்தில்}} — \emph{piṟkālattil} — \textquotedblleft{}in later times\textquotedblright{}.

\begin{quote}
\ta{பிற்காலத்தில்} \ta{அந்த} \ta{ஊர்} \ta{ஒரு} \ta{பெரிய} \ta{நகரமானது}.
\end{quote}

\begin{quote}
\emph{piṟkālattil anta ūr oru periya nakaramāṉatu.}
\end{quote}

\begin{quote}
In later times that village became a big city.
\end{quote}

\ta{பின்}, later + \ta{காலம்} + -\ta{இல்}; \ta{ன்} turns to \ta{ற்} before \ta{க}.

\begin{cousinweb}
\textbf{\ta{பிற்காலத்தில்}} (\emph{piṟkālattil}) joins native \ta{பின்}, \textquotedblleft{}what comes after\textquotedblright{} (DEDR 4205), with \ta{காலம்}, \textquotedblleft{}time\textquotedblright{}, from Sanskrit \emph{kāla} (CDIAL 3084), and -\ta{இல்}.
\end{cousinweb}

\begin{grammarlens}[title={What you've built}]
One more word for telling what happened.
\end{grammarlens}

\subsection*{Your turn}
\begin{itemize}
\raggedright
  \item \emph{Say it:} \emph{piṟkālattil}
  \item \emph{Say it:} \emph{piṟkālattil}, once more
  \item \emph{Say it:} \emph{piṟkālattil anta ūr oru periya nakaramāṉatu.} — In later times that village became a big city.
  \item \emph{Recall:} say the Tamil for to be staged, then the Tamil for a climax, then say \emph{piṟkālattil} again
\end{itemize}

\begin{tcolorbox}[breakable,colback=teal!4,colframe=teal!35!black,title={Before you move on}]
What does \ta{பிற்காலத்தில்} mean? (\textquotedblleft{}In later times\textquotedblright{}.) Say it once more.
\end{tcolorbox}
\section[\texorpdfstring{\ta{என்றென்றும்} (eṉṟeṉṟum)}{என்றென்றும் (eṉṟeṉṟum)}]{\ta{என்றென்றும்} (eṉṟeṉṟum) — forever; for all time}

\label{lesson:TA-C346-enrenrum}



\begin{quote}
Before the new one: say the Tamil for a climax, then the Tamil for in later times.
\end{quote}

\subsection*{You'll want to know: \ta{என்றென்றும்}}
\textbf{\ta{என்றென்றும்}} — \emph{eṉṟeṉṟum} — \textquotedblleft{}forever; for all time\textquotedblright{}.

\begin{quote}
\ta{அதன்பின்} \ta{அவர்கள்} \ta{என்றென்றும்} \ta{மகிழ்ச்சியாக} \ta{வாழ்ந்தார்கள்}.
\end{quote}

\begin{quote}
\emph{ataṉpiṉ avarkaḷ eṉṟeṉṟum makiḻcciyāka vāḻntārkaḷ.}
\end{quote}

\begin{quote}
After that they lived happily for ever.
\end{quote}

\ta{என்று}, when, doubled with -\ta{உம்}, every: at every time; a common close to a story.

\begin{cousinweb}
\textbf{\ta{என்றென்றும்}} (\emph{eṉṟeṉṟum}) is native Dravidian: DEDR 5151 lists \ta{என்று}, \textquotedblleft{}when\textquotedblright{}, and \ta{என்றும்}, \textquotedblleft{}always\textquotedblright{}, on the question base \ta{எ}-; here \ta{என்று} is doubled before -\ta{உம்}.
\end{cousinweb}

\begin{grammarlens}[title={What you've built}]
One more word for telling what happened.
\end{grammarlens}

\subsection*{Your turn}
\begin{itemize}
\raggedright
  \item \emph{Say it:} \emph{eṉṟeṉṟum}
  \item \emph{Say it:} \emph{eṉṟeṉṟum}, once more
  \item \emph{Say it:} \emph{ataṉpiṉ avarkaḷ eṉṟeṉṟum makiḻcciyāka vāḻntārkaḷ.} — After that they lived happily for ever.
  \item \emph{Recall:} say the Tamil for a climax, then the Tamil for in later times, then say \emph{eṉṟeṉṟum} again
\end{itemize}

\begin{tcolorbox}[breakable,colback=teal!4,colframe=teal!35!black,title={Before you move on}]
What does \ta{என்றென்றும்} mean? (\textquotedblleft{}Forever; for all time\textquotedblright{}.) Say it once more.
\end{tcolorbox}
\section[\texorpdfstring{\ta{இவ்வாறு} (ivvāṟu)}{இவ்வாறு (ivvāṟu)}]{\ta{இவ்வாறு} (ivvāṟu) — in this way}

\label{lesson:TA-C346-ivvaru}



\begin{quote}
Before the new one: say the Tamil for in later times, then the Tamil for forever.
\end{quote}

\subsection*{You'll want to know: \ta{இவ்வாறு}}
\textbf{\ta{இவ்வாறு}} — \emph{ivvāṟu} — \textquotedblleft{}in this way\textquotedblright{}.

\begin{quote}
\ta{இவ்வாறு} \ta{கதை} \ta{முடிந்தது}.
\end{quote}

\begin{quote}
\emph{ivvāṟu katai muṭintatu.}
\end{quote}

\begin{quote}
In this way the story ended.
\end{quote}

\ta{இ}, this + \ta{ஆறு}, way; \ta{அவ்வாறு}, in that way.

\begin{cousinweb}
\textbf{\ta{இவ்வாறு}} (\emph{ivvāṟu}) is built from the near-pointing \ta{இ}-, \textquotedblleft{}this\textquotedblright{}, and native \ta{ஆறு}, \textquotedblleft{}way, manner\textquotedblright{} (DEDR 405).
\end{cousinweb}

\begin{grammarlens}[title={What you've built}]
One more word for telling what happened.
\end{grammarlens}

\subsection*{Your turn}
\begin{itemize}
\raggedright
  \item \emph{Say it:} \emph{ivvāṟu}
  \item \emph{Say it:} \emph{ivvāṟu}, once more
  \item \emph{Say it:} \emph{ivvāṟu katai muṭintatu.} — In this way the story ended.
  \item \emph{Recall:} say the Tamil for in later times, then the Tamil for forever, then say \emph{ivvāṟu} again
\end{itemize}

\begin{tcolorbox}[breakable,colback=teal!4,colframe=teal!35!black,title={Before you move on}]
What does \ta{இவ்வாறு} mean? (\textquotedblleft{}In this way\textquotedblright{}.) Say it once more.
\end{tcolorbox}
\section[\texorpdfstring{\ta{விடி} (viṭi)}{விடி (viṭi)}]{\ta{விடி} (viṭi) — to grow light; to dawn}

\label{lesson:TA-C346-viti-dawn}



\begin{quote}
Before the new one: say the Tamil for forever, then the Tamil for in this way.
\end{quote}

\subsection*{You'll want to know: \ta{விடி}}
\textbf{\ta{விடி}} — \emph{viṭi} — \textquotedblleft{}to grow light; to dawn\textquotedblright{}.

\begin{quote}
\ta{கதை} \ta{முடிந்தபோது} \ta{பொழுது} \ta{விடிந்தது}.
\end{quote}

\begin{quote}
\emph{katai muṭintapōtu poḻutu viṭintatu.}
\end{quote}

\begin{quote}
By the time the story ended, day had dawned.
\end{quote}

Past (it …): \ta{விடிந்தது}, \emph{viṭintatu}.

Of the day: \ta{பொழுது} \ta{விடிந்தது}, it grew light.

\begin{cousinweb}
\textbf{\ta{விடி}} (\emph{viṭi}) is native Dravidian (DEDR 5475): to dawn, to break, as the day.
\end{cousinweb}

\begin{grammarlens}[title={What you've built}]
One more word for telling what happened.
\end{grammarlens}

\subsection*{Your turn}
\begin{itemize}
\raggedright
  \item \emph{Say it:} \emph{viṭi}
  \item \emph{Say it:} \emph{viṭi}, once more
  \item \emph{Say it:} \emph{katai muṭintapōtu poḻutu viṭintatu.} — By the time the story ended, day had dawned.
  \item \emph{Recall:} say the Tamil for forever, then the Tamil for in this way, then say \emph{viṭi} again
\end{itemize}

\begin{tcolorbox}[breakable,colback=teal!4,colframe=teal!35!black,title={Before you move on}]
What does \ta{விடி} mean? (\textquotedblleft{}To grow light; to dawn\textquotedblright{}.) Say it once more.
\end{tcolorbox}
\section[(dialogue)]{Looking back: legends, tales and how they end}

\label{lesson:TA-R346-review-340-346}



\begin{quote}
Say the Tamil for to decline and to renounce.
\end{quote}

\subsection*{Your turn}
Say these aloud:

\begin{itemize}
\raggedright
  \item to decline, to renounce, to exile, a golden age, a civilization
  \item an excavation, an antiquity, a treasure, to dig out, to come upon
  \item to restore, to engrave, ancient, an age, archaeology
  \item a legend, a myth, a made-up story, to make up, a proverb
  \item a monster, a fairy, a sorcerer, to curse, to change shape
  \item a storyteller, a moral tale, a bow-song, a setting, to be staged
  \item a climax, in later times, forever, in this way, to grow light
\end{itemize}

\begin{tcolorbox}[breakable,colback=teal!4,colframe=teal!35!black,title={Before you move on}]
To decline? (\textbf{\ta{வீழ்ச்சியடை}}, \emph{vīḻcciyaṭai}.) To renounce? (\textbf{\ta{துற}}, \emph{tuṟa}.) To exile? (\textbf{\ta{நாடுகடத்து}}, \emph{nāṭukaṭattu}.) A golden age? (\textbf{\ta{பொற்காலம்}}, \emph{poṟkālam}.) A civilization? (\textbf{\ta{நாகரிகம்}}, \emph{nākarikam}.) An excavation? (\textbf{\ta{அகழ்வாராய்ச்சி}}, \emph{akaḻvārāycci}.) An antiquity? (\textbf{\ta{தொல்பொருள்}}, \emph{tolporuḷ}.) A treasure? (\textbf{\ta{புதையல்}}, \emph{putaiyal}.) To dig out? (\textbf{\ta{தோண்டியெடு}}, \emph{tōṇṭiyeṭu}.) To come upon? (\textbf{\ta{கண்டெடு}}, \emph{kaṇṭeṭu}.) To restore? (\textbf{\ta{புனரமை}}, \emph{puṉaramai}.) To engrave? (\textbf{\ta{பொறி}}, \emph{poṟi}.) Ancient? (\textbf{\ta{பண்டைய}}, \emph{paṇṭaiya}.) An age? (\textbf{\ta{யுகம்}}, \emph{yukam}.) Archaeology? (\textbf{\ta{தொல்லியல்}}, \emph{tolliyal}.) A legend? (\textbf{\ta{பழங்கதை}}, \emph{paḻaṅkatai}.) A myth? (\textbf{\ta{புராணம்}}, \emph{purāṇam}.) A made-up story? (\textbf{\ta{கட்டுக்கதை}}, \emph{kaṭṭukkatai}.) To make up? (\textbf{\ta{புனை}}, \emph{puṉai}.) A proverb? (\textbf{\ta{பழமொழி}}, \emph{paḻamoḻi}.) A monster? (\textbf{\ta{அரக்கன்}}, \emph{arakkaṉ}.) A fairy? (\textbf{\ta{தேவதை}}, \emph{tēvatai}.) A sorcerer? (\textbf{\ta{மந்திரவாதி}}, \emph{mantiravāti}.) To curse? (\textbf{\ta{சபி}}, \emph{capi}.) To change shape? (\textbf{\ta{உருமாறு}}, \emph{urumāṟu}.) A storyteller? (\textbf{\ta{கதைசொல்லி}}, \emph{kataicolli}.) A moral tale? (\textbf{\ta{நீதிக்கதை}}, \emph{nītikkatai}.) A bow-song? (\textbf{\ta{வில்லுப்பாட்டு}}, \emph{villuppāṭṭu}.) A setting? (\textbf{\ta{கதைக்களம்}}, \emph{kataikkaḷam}.) To be staged? (\textbf{\ta{அரங்கேறு}}, \emph{araṅkēṟu}.) A climax? (\textbf{\ta{உச்சகட்டம்}}, \emph{uccakaṭṭam}.) In later times? (\textbf{\ta{பிற்காலத்தில்}}, \emph{piṟkālattil}.) Forever? (\textbf{\ta{என்றென்றும்}}, \emph{eṉṟeṉṟum}.) In this way? (\textbf{\ta{இவ்வாறு}}, \emph{ivvāṟu}.) To grow light? (\textbf{\ta{விடி}}, \emph{viṭi}.)
\end{tcolorbox}
\section[(dialogue)]{How It Ended — once more}

\label{lesson:TA-R346-review-final-346}



\begin{quote}
Say the Tamil for a climax and in later times.
\end{quote}

\subsection*{Your turn}
\begin{itemize}
\raggedright
  \item \emph{Say it:} a climax, in later times, forever, in this way, to grow light
\end{itemize}

\begin{tcolorbox}[breakable,colback=teal!4,colframe=teal!35!black,title={Before you move on}]
A climax? (\textbf{\ta{உச்சகட்டம்}}, \emph{uccakaṭṭam}.) In later times? (\textbf{\ta{பிற்காலத்தில்}}, \emph{piṟkālattil}.) Forever? (\textbf{\ta{என்றென்றும்}}, \emph{eṉṟeṉṟum}.) In this way? (\textbf{\ta{இவ்வாறு}}, \emph{ivvāṟu}.) To grow light? (\textbf{\ta{விடி}}, \emph{viṭi}.)
\end{tcolorbox}
